\setrunninghead{前言}
\section*{前言}
\addcontentsline{toc}{section}{\protect\textbf{前言}}
在开始这本书之前，我想给大家讲两个国家的故事。第一个国家，叫瑙鲁。第二个国家，叫挪威。

瑙鲁，是全球面积第三小的国家。这个国家曾经被称为“天堂岛”，因为岛上有非常丰富的磷矿。

当时的瑙鲁人，几乎不用上班，医疗、教育全部免费，房子、水电不要钱。许多瑙鲁人由于生活得太好，身患肥胖症。然而瑙鲁人并没有把自己的收入好好打理，而是随意地花掉了。

瑙鲁人原本以为磷矿可以开采150年，结果20年就开采完了。当磷矿开采完之后，瑙鲁就陷入了困局。最后，瑙鲁人不得不依靠变卖资产来维持生活。

同样是资源型的国家，挪威就做得好多了。

挪威是一个能源资源非常丰富的国家，但是人口比较少，产业也比较单一，主要依靠开采石油等能源来获取收入。如果石油等能源开采完了，那挪威会遇到很大的问题。

为了避免出现困境，挪威从1990年开始，设立了“挪威政府养老基金”。这只基金的目的，是储存石油收入并实现长期增长，最后实现财富的代际转移，即把当代人的石油收入，定期转化为更长期的金融产品收入。

这其实就是一个以国家为维度的超长期的定投。

挪威的投资策略，早期是配置40\%的股票、60\%的债券。其中，股票的参考基准是富时全球指数（FTSE All-World Index），这个指数是投资全球主要股票市场的指数；债券的参考基准，则是全球债券指数。

这样，挪威就把本土的石油收入，定期地转化为全球的“股票资产+债券资产”收入，分享全球股票和债券的收益。从此，挪威不用再担心自己发展掉队的问题了。

这保证了挪威的长期发展速度跟全球发展的平均速度相匹配。即便将来挪威跟当年的瑙鲁一样出现了资源枯竭，挪威人依靠自己这些年买入并积累下来的资产也可以生活得很好。

从2008年到2018年，这只基金获得了大约8\%的年化收益率，表现还是不错的。目前这只基金的规模接近1万亿美元，是全球最大的政府养老基金之一。

挪威人口只有529万，这相当于每个挪威人分配了18.9万美元的股票和债券的养老基金，未来还会继续增长。

通过这种方式，挪威实现了非常稳健的财富积累。在可以预见的未来，挪威的增长速度不会弱于全球平均增长速度。

这两个国家，其实就是我们日常生活中经常看到的两种不同的家庭。

有的家庭没有投资理财的意识，赚多少花多少。家庭成员一旦退休，就只有用大幅减少的退休金、养老金来维持生活，生活质量堪忧。而有的家庭，一开始就做好了计划，用收入定期买入优质的资产。当手里积累了足够多的资产之后，即便原来的收入减少或者消失，他们也可以依靠资产继续过美好的生活。

其实，无论是国家还是个人，在投资理财上，都是有共性的。

对挪威来说，是将石油收入转化为金融资产收入；对我们普通投资者来说，则是将工资收入转化为资产收入。石油有开采完的一天，退休后工资收入也会下降，但资产是可以长期为我们贡献源源不断的收入的。

那什么样的资产，适合用来长期定投？我们该如何通过定投获得更好的收益？又该如何避免投资中的风险？

这就是这本书要介绍给大家的。

第1章，介绍了定投前的准备，我们要具备哪些理念才能让我们的定投事半功倍。

第2章，介绍了什么是定投，其实很多人已经在定投了，只不过自己没有意识到。

第3章，介绍了如何通过定投来制订自己的财务自由计划，需要积累多少资产，才能让家庭过上理想的生活。

第4章，介绍了普通投资者必须掌握的投资品种。这些投资品种是我们手里的武器，灵活用好它们，我们能更快速地积累财富。

第5～7章，介绍了基金投资时最难的一类品种——股票指数基金。这类品种是长期定投收益最高的一类，也是风险很高的一类。

第8～9章，介绍了投资指数基金的高级技巧，包括如何挑选指数基金以及如何判断一只指数基金是否适合定投。

第10章，介绍了目前市面上最常见的几种定投技巧。我们可以付出更少的成本，获得更高的收益。

第11章，介绍了定投的卖出。俗话说，会买的是徒弟，会卖的是师傅。明白如何卖出，可以帮助我们落袋为安。

第12章，综合运用前面介绍的知识，帮助我们制订一个完整的定投计划，开启财务自由之旅。

第13章，介绍了投资者刚开始接触定投时最常遇到的10个问题。

我在创作这本书的过程中，得到了很多朋友的帮助。在创作这本书之初，我思考过一个问题：市面上讲定投的书很多，我来写定投的书，能给大家带来什么帮助呢？

后来想了一下，这几年我在网上分享了很多基金投资的内容，也得到了超过百万读者的关注。针对定投，大家最关心的内容是可以统计出来的，这其实就是大家对定投真正的疑惑。解答清楚这些疑惑，这本书也就水到渠成了。所以从某种意义上来说，这也是百万定投者一起创作出来的一本书，是属于定投者的一本书。因此，我创作这本书时最想感谢的人就是长期关注并支持螺丝钉的朋友们。你们的鼓励，是我坚持的最大动力。

也感谢出版社老师，她们在这本书编辑过程中不辞辛劳。中信出版社专业的编辑老师，是这本书坚实的后盾。

特别感谢我的妻子，这几年给了我很大的帮助，没有她的理解和支持，我很难走到今天。

读完这本书，其实只是一个小小的开始，大家可以关注我的微信公众号“定投十年赚十倍”，我会在每个交易日发布基金的估值数据以及关于基金投资的心得和感悟，也会组织一些线下活动，发放一些小福利。后续大家对书里的内容有任何疑问，或者投资中遇到任何问题，都欢迎与我讨论，我们互相学习，共同进步。
